
\section{Empirical Discovery: Mean-Dominated Geometry}
\label{sec:discovery}

\subsection{Activation Mean Is Statistically Significant}
\label{subsec:mean_significance}

We measure the noise-normalized mean magnitude
\[
R_x = \frac{\|\mu_x\|_2}{\sqrt{\mathbb{E}\|x\|_2^2 / N}}
\]
and show that activation means satisfy $R_x \gg 1$, while gradient-output means
satisfy $R_\delta \approx 1$.
This establishes that bias is intrinsic to activations and not to gradient noise.


\subsection{Outliers Are Dominated by the Mean}
\label{subsec:outlier_dominance}

We decompose activation tensors into mean, spike, and tail components using SVD,
and analyze the top $0.1\%$ largest-magnitude elements.
We show that the majority of squared magnitude within outliers is contributed
by the mean component.


\subsection{Spectral Role of Bias}
\label{subsec:spectral_role}

We analyze how the mean contributes to the spectral spike and quantify
its fraction of spike energy.
This reveals that the bias forms a dense, coherent direction embedded
within the dominant spectral subspace.


\subsection{Bias Propagation into Gradients}
\label{subsec:bias_propagation}

Using the gradient identity
\[
\mathbb{E}[\delta x^\top]
=
\mathbb{E}[\delta]\mu_x^\top
+
\mathbb{E}[\delta \tilde{x}^\top],
\]
we show that activation bias induces a coherent rank-one component
in gradient matrices even when gradient outputs themselves have negligible mean.


\section{The Blessing: Efficient Mean Removal}
\label{sec:mean_removal}

\subsection{Mean Removal as Rank-One Suppression}

We define
\[
X_c = X - \mathbf{1}\mu^\top
\]
and show that mean removal significantly reduces extreme outliers and
partially attenuates spectral spike energy.


\subsection{Element-Wise and Block-Wise Outlier Reduction}

We quantify p99.9 entry magnitude and blockwise maxima before and after
mean removal, demonstrating substantial reduction in extreme scales.


\section{Downstream Validation}
\label{sec:applications}

\subsection{Low-Bit Quantization}
\label{subsec:quantization}

We analyze how mean removal improves signal-to-noise ratio under
blockwise FP4 quantization and compare against SVD-based spike whitening.


\subsection{Optimizer Conditioning}
\label{subsec:optimizer}

We examine the effect of bias on Adam second-moment estimates
\[
v = \mathbb{E}[g \odot g]
\]
and show that mean removal reduces pathological inflation of adaptive
learning-rate scaling.


\section{Theoretical Formalization}
\label{sec:theory}

We formalize a minimal model of activation bias and semantic variability,
derive bounds on mean-dominated outlier contributions, and prove that
mean removal preserves semantic covariance up to controlled error.


\section{Discussion and Implications}
\label{sec:discussion}

We discuss:
\begin{itemize}
    \item Why bias emerges in language representations.
    \item When mean removal suffices versus when spectral whitening is needed.
    \item Broader implications for scaling, stability, and hardware efficiency.
\end{itemize}


\section{Conclusion}
\label{sec:conclusion}

We argue that mean-dominated bias is a fundamental structural component
of LLM representations and a key hidden driver of instability.
By recognizing bias as both curse and blessing, we obtain a simple,
efficient, and theoretically grounded mechanism for stabilizing
large-scale training.

\newpage


\bibliographystyle{../doc/iclr-template/master/iclr2026/iclr2026_conference}
\bibliography{references}

\newpage

\appendix
%\input{proofs}

